\section{Introduction}
\label{sec:intro}

Simulation-based testing is vital for verifying safety-critical Autonomous Cyber-Physical Systems (ACPS), ranging from Autonomous Driving Systems (ADS) to Unmanned Aerial Vehicles (UAVs)~\cite{sdcscissor,sensodat}.
High-fidelity physics simulators (e.g., BeamNG.tech and Gazebo/PX4) allow engineers to evaluate control policies across hazardous corner cases without risking physical damage.
However, simulating multi-agent interactions, sensor telemetry, and nonlinear friction forces is computationally prohibitive.
A realistic verification suite containing thousands of test scenarios requires hundreds of machine-hours to execute.
To mitigate this bottleneck, \emph{test case prioritization} (TCP)~\cite{apfd,yooharman} reorders test suites so that failure-revealing scenarios are executed as early as possible.
Prioritization performance is standardly evaluated via the Average Percentage of Faults Detected (APFD)~\cite{apfd}:
\begin{equation}
\mathrm{APFD}(\pi) = 1 - \frac{\sum_{i=1}^{m} TF_i}{n \cdot m} + \frac{1}{2n},
\label{eq:apfd}
\end{equation}
where $n$ is the total number of test cases, $m$ is the count of failing test cases, and $TF_i$ is the rank position of the first test exposing fault $i$.
An ideal prioritization yields an APFD approaching $1.0$, whereas random execution yields $\approx 0.50$.

Existing approaches to autonomous system test prioritization broadly fall into three categories:
(i)~\emph{Heuristic and search-based diversification} (e.g., SDC-Prioritizer~\cite{sdcprio} and Surrealist~\cite{roadfury}), which employ genetic algorithms to maximize behavioral or static diversity;
(ii)~\emph{Feature-based machine learning} (e.g., SDC-Scissor~\cite{sdcscissor,saner}), which extracts hand-crafted tabular descriptors (such as turn counts, bounding box areas, and mean curvatures); and
(iii)~\emph{Deep sequence models} (e.g., ITS4SDC~\cite{its4sdc} and RoadFury~\cite{roadfury}), which process waypoint sequences using LSTMs or Transformers.
While deep sequence architectures achieve strong APFD ($\approx 0.807$), a critical examination reveals four foundational limitations:
\begin{enumerate}
    \item \textbf{Coordinate Frame Fragility}: Existing neural prioritizers rely on extrinsic coordinate features (Cartesian coordinates $(x, y, z)$ or yaw angles). Consequently, rotating a scenario's global coordinate frame induces severe rank drift ($\Delta\mathrm{APFD} > 0.05$), even though the underlying physical trajectory is identical.
    \item \textbf{Discretization Sensitivity}: Simulators frequently alter trajectory waypoint sampling frequencies. Fixed-rate neural networks suffer performance degradation when evaluated on resampled curves.
    \item \textbf{Unphysical Failure Predictions}: Standard statistical loss functions treat risk prediction as unconstrained pattern matching, frequently producing non-monotonic predictions where sharper curves or extreme aerodynamic gusts are paradoxically predicted as less hazardous.
    \item \textbf{Lack of Statistical Guarantees}: Existing prioritizers provide heuristic rankings with zero finite-sample safety assurances, leaving test engineers without rigorous stopping criteria.
\end{enumerate}

To resolve these limitations, we propose \textbf{CliffordTrajNet}, a unified geometric deep learning and continuous state-space framework for test prioritization across autonomous vehicles and UAVs.
First, CliffordTrajNet formulates spatial trajectories in Clifford Geometric Algebra $C\ell(p, q)$ ($C\ell(2,0)$ for 2D road driving and $C\ell(3,0)$ for 3D aerial flight).
By operating on grade-preserving multivectors and rotor sandwich operations ($v' = R v R^\dagger$), CliffordTrajNet guarantees exact planar $SE(2)$ invariance and substantially reduces 3D spatial sensitivity without gimbal lock.
Second, we incorporate a continuous-time Selective State Space Model (\textbf{TrajMamba}) driven by arclength ODE integration ($\Delta t = \Delta s / v$), substantially reducing discretization sensitivity across variable waypoint densities.
Third, we formulate a physics-informed monotonicity regularizer enforcing curvature monotonicity, reducing unphysical predictions by $5.6\times$.
Finally, we integrate \textbf{Monotone Conformal Risk Control} (CRC), providing finite-sample probably approximately correct (PAC) statistical bounds that guarantee test failure miss-rates remain bounded under tight execution budgets.

In summary, this paper presents the following contributions:
\begin{itemize}
    \item \textbf{Invariant Clifford Geometric Foundations}: We formulate test prioritization in Clifford Geometric Algebra, guaranteeing exact $SE(2)$ invariance on ground roads ($\Delta\mathrm{APFD} = 0.0000$) and high spatial robustness on 3D UAV paths ($\Delta\mathrm{APFD} \le 0.0119$).
    \item \textbf{Continuous State-Space Dynamics (TrajMamba)}: We achieve strong robustness against waypoint discretization sensitivity ($\Delta\mathrm{APFD} \le 0.0012$) via continuous arclength ODE integration over variable-resolution trajectories.
    \item \textbf{PAC Safety Bounds}: We establish rigorous stopping criteria, showing that allocating just $17.05\%$ budget catches $\ge 98.72\%$ of critical simulation failures on evaluated benchmarks.
    \item \textbf{Extensive Cross-Domain Evaluation}: Across 32,580 SensoDat simulation tests, nine external SDC benchmarks, and the official SBFT 2026 UAV benchmark, CliffordTrajNet achieves the highest AUC among evaluated methods ($0.9399$) while matching the strongest APFD performance ($0.8063$ vs. $0.8066$ for RoadFury), outperforming heuristic, tabular, and deep learning baselines.
\end{itemize}
